\scieloAbstractContainer{RESUMO}{Foi feito um estudo sobre a estrutura diamétrica e o arranjo espacial das espécies mais abundantes
que ocorrem em um dos cinco fragmentos florestais da Fazenda Experimental Edgardia (1.200  ha),
município de Botucatu, SP. A área do remanescente é de 56 ha. Inventariaram-se 185 unidades
amostrais contíguas de 100 m 2 , que compuseram uma amostra de 1,85 ha. Levantou-se a
composição florística e estrutura horizontal da vegetação do centro do remanescente. Dentre as
dez espécies de maior valor de importância, sete apresentaram abundância que permitiu realizar
os estudos sobre distribuição diamétrica e distribuição
espacial. O somatório dos índices de valores de importância (IVI) obtidos para as espécies mais abundantes (185,98) mostrou que elas
têm uma grande participação na estrutura da vegetação que está se regenerando. Após ajuste das
distribuições diamétricas das sete espécies, seis apresentaram um padrão de distribuição na forma
de J invertido, refletindo regeneração contínua dos indivíduos dessas espécies. Quatro dessas
distribuições foram distinguidas como tipo I (ou distribuição balanceada). As características
relacionadas com as distribuições de diâmetros ajustadas auxiliaram na indicação de uma
atenuação da perturbação antropogênica. Verificou-se relação direta entre o arranjo espacial e
a síndrome de dispersão para Trichilia clausseni e Metrodorea nigra e, dependendo das escalas
analisadas, para T. catigua e T. casaretti.}{Palavras-chave:}{comunidades, distribuição diamétrica, estatística espacial, estrutura, florestas, florística, padrão}

\begin{multicols}{2}
\scieloSectionContainer{Introduction}
\par{}Brain interfaces are commonly considered as the
ultimate means of interaction between humans and
machines or computers (Allison et al., 2012; Doud et al.,
2011; Hochberg et al., 2012; Sirvent Blasco et al.,
2012). The idea of controlling our environment using
brain activity opens the possibility for many potential
applications using these interfaces. For instance, a

\scieloSectionContainer{4. CONCLUSÕES}

\lipsum

\scieloSectionContainer{Materials and Methods}

\scieloSubSectionContainer{Animals}

\par{}Specimens representing 12 species of polyclads were
collected from different sites in the Caribbean. In addition,
a deep-sea specimen of \textit{Anocellidus profundus} Quiroga,
Bolaños and Litvaitis, 2006, from the North Pacific was
also included in the analysis (specimen courtesy of Dr.
Janet Voight, Field Museum of Natural History, Chicago,
Illinois, USA). For details of collection, fixation, and spe-
cies identification see Litvaitis \textit{et al}. (2010).

\scieloSubSectionContainer{Histological studies of polyclad nervous systems}

\par{}A specimen of each species was embedded in paraffin
in its entirety, longitudinally sectioned at 7-10 mm, and
stained with Milligan Trichrome technique which allows
differentiating well among connective tissue, muscles, and
nervous fibers (Presnell and Schreibman, 1997). Additional
cross and sagittal sections were also prepared. De-paraffi-
nized sections were treated with 3\% potassium dichro-
mate-hydrochloric acid solution for 5 min. Following a
distilled water rinse, sections were stained in acid fuchsin
for 8 min. After a second distilled water rinse, they were
placed into 1\% phosphomolybdic acid for 2 min and then
stained with 2\% solution of Orange G for 5 min. After a fi-
nal rinse with distilled water, the sections were treated with
1\% hydrochloric acid solution for 2 min stained in Fast
Green for 8 min, treated with 1\% acetic acid for 3 min and
then rinsed in 95\% alcohol and dehydrated. Finally, the sec-
tions were cleared with Histoclear (National Diagnostics)
and mounted in Permount (Fisher Scientific). Slides were
observed and photographed under an Axiostar Plus (Zeiss,
Thornwood, New York) light microscope. We applied the
morphological criteria of Reuter et al. (1998) to distinguish
main nerve cords from secondary nerve cords for the exam-
ination of the NS. Cross and sagittal sections were used for
interpretation. In addition, whole mounts of all species
were prepared, by dehydrating the specimens in a graded
alcohol series, cleared with Histoclear, and mounted in
Permount.

\lipsum

\scieloImageContainerOneCol{0.5\textwidth}{img1.jpg}{meu titulo}

%\begin{center}
%{
%\vspace{8mm}
%\centerline{
%\includegraphics[width=\maxwidth{0.5\textwidth}]{img1.jpg}
%}
%\vspace{8mm}
%}
%\end{center}

\lipsum

\scieloImageContainerTwoCol{img2.jpg}{\textbf{Figure 1:} A specimen of each species was embedded in paraffin in its entirety, longitudinally sectioned}

\lipsum
\lipsum

\scieloImageContainerTwoCol{img2.jpg}{\textbf{Figure 1:} A specimen of each species was embedded in paraffin in its entirety, longitudinally sectioned}

\lipsum
\lipsum

Os resultados no presente trabalho mostraram que a espécie tem
um padrão de distribuição espacial agregado. Sendo
assim, a localização dos indivíduos de M. nigra pode
ser usada para identificar os locais menos perturbados
dentro do mosaico sucessional.

\scieloSectionContainer{4. CONCLUSÕES}

Com base nos resultados do presente trabalho,
concluiu-se que:

\begin{scieloReferencesContainer}[REFERÊNCIAS]

\scieloReferencesItem{Baddeley A. \textit{Analysing spatial point patterns in} R. Perth:
CSIRO:University of Western Australia; 2010. Workshop Notes, Version 4.1.}

\scieloReferencesItem{Carvalho WA, Panoso LA, Moraes MH. \textit{Levantamento
semidetalhado dos solos da fazenda experimental Edgardia
– município de Botucatu}. Botucatu: Faculdade de Ciências
Agronômicas – UNESP; 1991. Boletim Científico vol.
2. 467 p.}

\scieloReferencesItem{Cruz-Rot M. Introdución al análisis de datos mapeados
o algunas de las (muchas) cosas que puedo hacer si tengo
coordenadas. \textit{Ecosistemas} 2006; 15(3): 19-39.}

\scieloReferencesItem{King GJ, Keeling SD, Hocevar RA, Wheeler TT. The timing of treatment for Class II malocclusions in children: a literature review. Angle
Orthod. 1990;60(2):87-97. PMid:2111647.}

\scieloReferencesItem{Proffit WR. The timing of early treatment: an overview. Am J Orthod Dentofacial Orthop. 2006 Apr;129(4 Suppl):S47-9. http://dx.doi.
org/10.1016/j.ajodo.2005.09.014. PMid:16644417.}

\scieloReferencesItem{Franchi L, Baccetti T, De Toffol L, Polimeni A, Cozza P. Phases of the dentition for the assessment of skeletal maturity: a diagnostic performance
study. Am J Orthod Dentofacial Orthop. 2008 Mar;133(3):395-400. PMid: 18331939.}

\scieloReferencesItem{Bullock TH and Horridge GA (1965) Structure
and Function in the Nervous Systems of Invertebrates.
Vol. 2. W.H. Freeman and Co, San
Francisco, 1719 pp.}

\scieloReferencesItem{Carranza S, Baguñà J and Riutort M (1997)
Are the Platyhelminthes a monophyletic primitive
group? An assessment using 18S rDNA
sequences. Mol Biol Evol 14:485-497.}

\scieloReferencesItem{Cebrià F (2008) Organization of the nervous system in the model
planarian Schmidtea mediterranea: An immunocytochemical study. Neurosci Res 61:375-384.}

\scieloReferencesItem{Chien P and Koopowitz H (1972) The ultrastructure of neuromuscular systems in Notoplana acticola, a free-living polyclad flatworm. Z Zellforsch Mikroskop Anat 133:277-288.}


\scieloReferencesItem{Chien PK and Koopowitz H (1977) Ultrastructure of nerve plexus
in flatworms. III. The infra-epithelial nervous system. Cell
Tissue Res 176:335-347.}

\scieloReferencesItem{Day TA, Maule AG, Shaw C and Pax RA (1997) Structure-activity relationships of FMRFamide-related peptides contracting Schistosoma mansoni muscle. Peptides 18:917-921.}

\scieloReferencesItem{Egger B, Gschwentner R and Rieger R (2007) Free-living
flatworms under the knife: Past and present. Dev Genes Evol
217:89-104.}

\scieloReferencesItem{Egger B, Lapraz F, Tomiczek B, Müller S, Dessimoz C, Girstmair
J, Skunca N, Rawlinson KA, Cameron CB, Beli E et al.
(2015) A transcriptomic-phylogenomic analysis of the evolutionary relationships of flatworms. Curr Biol 25:1-7.}

\scieloReferencesItem{Ehlers U (1985) Das Phylogenetische System der Plathelminthes.
Gustav Fischer Verlag, Stuttgart, 317 pp.}

\scieloReferencesItem{Fernandes MC, Alvares EP, Gama P and Silveira M (2003) Serotonin in the nervous system of the head region of the land planarian Bipalium kewense. Tissue Cell 35:479-486.}

\scieloReferencesItem{Forest DL and Lindsay SM (2008) Observations of serotonin and
FMRFamide-like immunoreactivity in palp sensory structures and the anterior nervous system of spionid polychaetes. J Morphol 269:544-551.}

\scieloReferencesItem{Girstmair J, Schnegg R, Telford MJ and Egger B (2014) Cellular
dynamics during regeneration of the flatworm Monocelis sp. (Proseriata, Platyhelminthes). Evo Devo 5:e37.}

\scieloReferencesItem{Golding DM (1992) Polychaeta: Nervous system. In: Harrison
FW and Gardiner SL (eds) Microscopic Anatomy of Invertebrates. Vol. 7. Wiley-Liss, New York, pp 155-179.}

\scieloReferencesItem{Gustafsson MKS, Halton DW, Kreshchenko ND, Movsessian SO,
Raikova OI, Reuter M and Terenina NB (2002) Neuropeptides in flatworms. Peptides 23:2053-2061.}

\scieloReferencesItem{Hadenfeldt D (1929) Das Nervensystem von Stylochoplana
maculata und Notoplana atomata. Z wiss Zool 133:586638.}

\scieloReferencesItem{Halton DW and Gustafsson MKS (1996) Functional morphology
of the platyhelminth nervous system. Parasitology
113:S47-S72.}

\end{scieloReferencesContainer}

%\scieloLicenseContainer{License information: This is an open-access article distributed under the terms of the Creative Commons Attribution License, which permits unrestricted use, distribution, and reproduction in any medium, provided the original work is properly cited.}

\newpage

Os resultados no presente trabalho mostraram que a espécie tem
um padrão de distribuição espacial agregado. Sendo
assim, a localização dos indivíduos de M. nigra pode
ser usada para identificar os locais menos perturbados
dentro do mosaico sucessional.

\scieloSectionContainer{4. CONCLUSÕES}

Com base nos resultados do presente trabalho,
concluiu-se que:

Os resultados no presente trabalho mostraram que a espécie tem
um padrão de distribuição espacial agregado. Sendo
assim, a localização dos indivíduos de M. nigra pode
ser usada para identificar os locais menos perturbados
dentro do mosaico sucessional.

\scieloSectionContainer{5. CONCLUSÕES}

Com base nos resultados do presente trabalho,
concluiu-se que:


\end{multicols}
